\section{Mask dan Instruksi TEST}
\label{sec:mask-test}

Kemampuan operasi logika yang paling sering dipakai adalah menyentuh
satu bit tertentu tanpa mengganggu bit lain. Untuk itu diperlukan
\term{mask}, yaitu pola bit yang dirancang agar operasi logika
menghasilkan efek yang dikehendaki pada posisi yang dituju.

\subsection{Menyalakan, Memadamkan, Membalik, dan Menguji Bit}

Bila bit yang diincar berada pada posisi ke-$n$ (dihitung dari nol
pada bit paling kanan), mask-nya adalah $2^n$, yaitu nilai dengan
tepat satu bit bernilai 1 pada posisi tersebut. Misalkan \reg{AL}
berisi \hex{A5} = \bin{1010 0101} dan bit yang disasar adalah bit
ke-3 (\hex{08} = \bin{0000 1000}).

\begin{table}[H]
    \centering
    \caption{Pola manipulasi satu bit tertentu dengan mask.}
    \label{tab:manipulasi-bit}
    \small
    \begin{tabular}{@{}p{2.4cm}p{3.4cm}p{3.6cm}p{3.6cm}@{}}
        \toprule
        \textbf{Kehendak} & \textbf{Mask} & \textbf{Instruksi} &
        \textbf{Hasil pada \hex{A5}, bit 3} \\
        \midrule
        Menyalakan bit & mask & \lstasm|OR AL, mask| &
        \bin{1010 1101} = \hex{AD} \\
        \addlinespace
        Memadamkan bit & \lstasm|NOT mask| &
        \lstasm|AND AL, mask| & \bin{1010 0101} = \hex{A5} \\
        \addlinespace
        Membalik bit & mask & \lstasm|XOR AL, mask| &
        \bin{1010 1101} = \hex{AD} \\
        \addlinespace
        Menguji bit & mask & \lstasm|TEST AL, mask| &
        \flag{ZF} = 0 berarti bit 3 menyala \\
        \bottomrule
    \end{tabular}
\end{table}

Pola dasarnya:

\begin{lstlisting}[language=MASM, numbers=none]
OR  AL, 08h         ; nyalakan bit 3
AND AL, 0F7h        ; padamkan bit 3 (F7h = NOT 08h)
XOR AL, 08h         ; balikkan bit 3
TEST AL, 08h        ; uji bit 3 tanpa mengubah AL
JNZ BIT_MENYALA     ; lompat bila bit 3 bernilai 1
\end{lstlisting}

Mask penghapus \hex{F7} diperoleh dari \lstasm|NOT 08h|. Bila mask
dihitung saat program berjalan, mask disimpan ke register, dibalik
dengan \asm{NOT}, lalu dipakai pada \asm{AND}.

\subsection{TEST dan CMP}

\asm{TEST} dan \asm{CMP} sama-sama melakukan operasi tanpa menyimpan
hasil; yang tertinggal hanya jejak pada register \eng{flag}.
Perbedaannya terletak pada operasi yang dilakukan.

\begin{table}[H]
    \centering
    \caption{Perbandingan TEST dan CMP.}
    \label{tab:test-cmp}
    \small
    \begin{tabular}{@{}p{2.2cm}p{4.6cm}p{5.8cm}@{}}
        \toprule
        \textbf{Aspek} & \asm{TEST} & \asm{CMP} \\
        \midrule
        Operasi internal & \asm{AND} tanpa menyimpan hasil &
        \asm{SUB} tanpa menyimpan hasil \\
        \addlinespace
        \flag{CF} dan \flag{OF} & Selalu 0 &
        Sesuai hasil pengurangan \\
        \addlinespace
        Pemakaian utama & Memeriksa keadaan bit tertentu &
        Membandingkan dua nilai secara utuh \\
        \addlinespace
        Lompatan lazim & \asm{JZ}/\asm{JNZ}, \asm{JP}/\asm{JNP} &
        \asm{JE}/\asm{JNE}, \asm{JA}/\asm{JB} \\
        \addlinespace
        Contoh & \lstasm|TEST AL, 01h| lalu \asm{JNZ} &
        \lstasm|CMP AL, BL| lalu \asm{JA} \\
        \bottomrule
    \end{tabular}
\end{table}

Gunakan \asm{TEST} bila yang dipermasalahkan adalah bit tertentu, dan
\asm{CMP} bila yang dibandingkan adalah nilai keseluruhan.
Kebiasaan khusus \asm{TEST} adalah \lstasm|TEST AL, AL|, yang
melakukan \asm{AND} antara \reg{AL} dengan dirinya sendiri sehingga
isinya tidak berubah, tetapi \flag{ZF}, \flag{SF}, dan \flag{PF}
diperbarui.

\begin{lstlisting}[language=MASM, numbers=none]
TEST AL, 80h        ; periksa bit tanda
JNZ NEGATIF         ; ZF = 0 -> bit tanda menyala

TEST BL, 04h        ; periksa bit ke-2
JZ  BIT_NOL         ; ZF = 1 -> bit ke-2 padam

TEST AL, AL         ; periksa apakah AL = 0
JZ  SELESAI
\end{lstlisting}

\begin{peringatanbox}
\vspace{1.1em}
Kesalahan yang sering terjadi adalah menuliskan mask secara terbalik.
Untuk memadamkan bit ke-3, instruksi yang benar adalah
\lstasm|AND AL, 0F7h| (bit ke-3 pada mask bernilai 0), bukan
\lstasm|AND AL, 08h| yang justru menghapus seluruh bit lain.
\end{peringatanbox}
